\subsection{分次正则环}

设 \(A_{0}\) 为环，P 为一个 \(A_{0}\) 上的 N 型分次模且次数 \textgreater{} 0。用 A 表示环 \(\mathsf{S}_{\mathrm{A}_{0}}(\mathrm{P})\)；A 上存在唯一的 N-分次结构，使 \(A_{0}\) 的次数为 0 且 P 是 A 上的一个子 \(A_{0}\) -分次模。用 \(A_{+}\) 表示 A 的理想 \(\bigoplus_{n>0} A_{n}\)。则典范映射 \(P \longrightarrow A_{+}/A_{+}^{2}\) 是 \(A_{0}\) -分次模的同构 (参见 A, III, 第 76 页, 命题 10)。

若 \(A_{0}\) -模 P 是分次自由的 (A, II, 第 167 页, 注 3)，且 \((x_{i})_{i\in I}\) 是由齐次元素组成的 P 的一组基，则 \(A_{0}\) -分次代数 \(\mathbf{S}_{\Lambda_{0}}(\mathbf{P})\) 同构于多项式代数 \(\mathrm{A}_{0}[(\mathrm{X}_{i})_{i\in I}]\)，后者赋以使每个 \(X_{i}\) 为次数 \(\deg(x_{i})\) 的齐次元素的分次结构。我们称每个 N 型 \(A_{0}\) -分次代数为 \(A_{0}\) -分次多项式代数，如果它同构于前述形状的一个 \(A_{0}\) -分次代数。

当环 \(\Lambda_0\) 正则且 \(\mathrm{A}_0\) -模 P 是有限生成投射的时，环 \(\mathsf{S}_{\mathrm{A}_0}(\mathrm{P})\) 正则 ( \(\S 4\) , 见第 2 节, 命题 4 的推论 4)。反之：

定理 2.--- 设 A 为正则环，N 型分次的。由 A 中次数为 0 的元素组成的环 A₀ 是正则的；存在一个次数 \textgreater{} 0 的有限生成分次投射 A₀-模 P，使 A 作为分次 A₀-代数同构于 \(S_{A_{0}}(P)\)。

用 P 表示 \(A_{0}\) 分次模 \(A_{+}/A_{+}^{2}\)。由第 3 节的命题 4，环 \(A_{0}\) 正则且 \(A_{0}\) -模 P 是投射的且有限生成的。于是 P 的齐次分量都是投射的，且存在一个 \(A_{0}\) -线性截面 \(\varphi: P \to A_{+}\)（它是次数为 0 的分次映射，作用于典范满射 \(A_{+} \to P\)）。设 \(f: S_{\mathrm{A}_{0}}(\mathrm{P}) \longrightarrow \Lambda\) 为 \(A_{0}\) -代数同态，它扩张 \(\varphi\)。由命题 4，f 扩张为 \(\mathbf{S}_{\mathrm{A}_{0}}(\mathrm{P})\) 关于 \(\mathbf{S}_{\mathrm{A}_{0}}(\mathrm{P})_{+}\) -进拓扑的分离完备化到 A 关于 \(A_{+}\) -进拓扑的分离完备化上的一个同构。因此 f 是单射的，且它的像在 A 中对 \(A_{+}\) -进拓扑是稠密的。但由于 A 的齐次分量上的诱导拓扑是离散的，且 f 的像是一个分次子模，这就蕴含 f 是双射的。

推论 1. - 设 B 为正则环，带正次数的分次环。假设每个有限型投射 \(\mathbf{B}_0\) -模都是自由的。

\begin{enumerate}
\def\labelenumi{\alph{enumi})}
\item
  环 B \(_{0}\) 是整环且正则，且 B \(_{0}\) -代数 B 是一个 B \(_{0}\) 分次多项式代数，有限型的。
\item
  设 A 为 B 的一个分次子环使 \(A_{0} = B_{0}\)，且 B 是有限生成 A-模。下列条件等价：
  \begin{enumerate}
  \def\labelenumii{(\roman{enumii})}
  \item
    A-模 B 是分次自由的；
  \item
    A-模 B 是平坦的；
  \item
    A 是一个 \(A_{0}\) 分次多项式代数，有限型的。
  \end{enumerate}
\end{enumerate}

由定理 2，环 B₀ 正则，从而是正则整环的乘积 (§ 4, 第 2 节, 例 2)。对 B₀ 的每个幂等元 e，B₀-模 B₀e 是投射的，从而是自由的，这蕴含 e = 0 或 1；于是 B₀ 是整环。于是断言 a) 由定理 2 及有限型分次投射 B₀-模是分次自由的这一事实得出。

我们来证明 b)。

\begin{enumerate}
\def\labelenumi{(\roman{enumi})}
\item
  \(\Leftrightarrow\) (ii)：有限型分次平坦 \(A_0\) -模都是分次自由的 (II, § 5, 第 2 节, 定理 1 的推论 2)。于是 (i) 与 (ii) 的等价性由 A, X, 第 144 页, 命题 8 得出。
\item
  \(\Leftrightarrow\) (iii)：由于 B 在 A 上是整的且是有限生成 A \(_0\) -代数，A 是有限型 A \(_0\) -代数 (V, § 1, 第 9 节, 引理 5)，从而是 Noether 环。于是 (ii) 与 (iii) 的等价性由 § 4, 第 5 节的命题 8 的推论以及应用到 A 的 a) 得出。
\end{enumerate}

推论 2.- 设 \(k\) 为域，B 为一个 \(k\) -分次多项式代数，有限型的，A 为 B 的一个分次子代数。下列条件等价：

\begin{enumerate}
\def\labelenumi{(\roman{enumi})}
\item
  B 是分次自由 Λ-模；
\item
  B 是平坦 A-模；
\item
  有 \(\mathrm{Tor}_1^\Lambda (k,\mathbf{B}) = 0\)；
\item
  代数 A 是一个有限生成分次多项式 k-代数，且 A 的每个由齐次元素组成的代数自由生成序列都是 B-正则的。
\end{enumerate}

蕴含关系 (i) \(\Rightarrow\) (ii) 与 (ii) \(\Rightarrow\) (iii) 是显然的，而蕴含关系 (iii) \(\Rightarrow\) (i) 由 \(\Lambda\) , X, 第 144 页, 命题 8 a) 得出。

(i)⇒(iv)：由于环 B 正则且在 A 上忠实平坦，环 A 由 I, § 3, 第 6 节的命题 11 是 Noether 的，且由 § 4, 第 5 节的命题 8 是正则的，从而是一个分次多项式 k-代数 (推论 1, a))。A 的每个代数自由生成序列都是 A-正则的 (A, X, 第 158 页, 例)，从而由于 B 在 A 上平坦而是 B-正则的。

((iv) \(\Rightarrow\) (iii))：设条件 (iv) 满足，并设 x 为 A 的一个由齐次元素组成的代数自由生成序列。由于序列 x 是 A-正则的，它对 A 是完全割的 (A, X, 第 157 页, 命题 5)，于是 A-模 \(\mathrm{Tor}_{1}^{\mathrm{A}}(k,\mathrm{B})\) 同构于 \(\mathrm{H}_{1}(\mathbf{x},\mathrm{B})\) (A, X, 第 159 页, 注 3)；但后者为零，因为序列 x 是 B-正则的。

推论 1 与推论 2 蕴含 LIE, V, § 5, 第 5 节的引理 5。

